% TeX program = pdflatex
\documentclass[11pt,a4paper]{amsart}

\newcommand{\notelanguage}{finnish}
% Shared preamble for course notes and exercise sets.

\usepackage[T1]{fontenc}
\usepackage{lmodern}
\providecommand{\notelanguage}{english}
\usepackage[finnish,english]{babel}
\AtBeginDocument{\expandafter\selectlanguage\expandafter{\notelanguage}}

\newcommand{\notesfinnishlanguage}{finnish}
\ifx\notelanguage\notesfinnishlanguage
\newcommand{\notestheoremname}{Lause}
\newcommand{\noteslemmaname}{Lemma}
\newcommand{\notespropositionname}{Propositio}
\newcommand{\notescorollaryname}{Seurauslause}
\newcommand{\notesdefinitionname}{Määritelmä}
\newcommand{\notesexamplename}{Esimerkki}
\newcommand{\notesproblemname}{Tehtävä}
\newcommand{\notesexercisename}{Tehtävä}
\newcommand{\notesremarkname}{Huomautus}
\newcommand{\notessolutionname}{Ratkaisu}
\else
\newcommand{\notestheoremname}{Theorem}
\newcommand{\noteslemmaname}{Lemma}
\newcommand{\notespropositionname}{Proposition}
\newcommand{\notescorollaryname}{Corollary}
\newcommand{\notesdefinitionname}{Definition}
\newcommand{\notesexamplename}{Example}
\newcommand{\notesproblemname}{Problem}
\newcommand{\notesexercisename}{Exercise}
\newcommand{\notesremarkname}{Remark}
\newcommand{\notessolutionname}{Solution}
\fi

\usepackage[a4paper,margin=30mm]{geometry}
\usepackage{microtype}
\usepackage{graphicx}
\usepackage{enumitem}

\newcommand{\exerciseheading}[2]{%
  \par\noindent
  {\large\bfseries #1\par}
  \vspace{0.1em}
  \noindent{\large #2\par}
  \vspace{1.5em}
}

\usepackage{amsmath}
\usepackage{amssymb}
\usepackage{amsthm}
\usepackage{mathtools}
\usepackage{aliascnt}

% Linear systems: the argument gives the number of alignment pairs.
\newenvironment{linearsystem}[1]
  {\left\{\begin{alignedat}{#1}}
  {\end{alignedat}\right.}

\usepackage[hidelinks,pdfusetitle]{hyperref}

\numberwithin{equation}{section}
\setlist{itemsep=0.25em,topsep=0.5em}

\theoremstyle{plain}
\newtheorem{theorem}{\notestheoremname}[section]
\newaliascnt{lemma}{theorem}
\newtheorem{lemma}[lemma]{\noteslemmaname}
\aliascntresetthe{lemma}
\newaliascnt{proposition}{theorem}
\newtheorem{proposition}[proposition]{\notespropositionname}
\aliascntresetthe{proposition}
\newaliascnt{corollary}{theorem}
\newtheorem{corollary}[corollary]{\notescorollaryname}
\aliascntresetthe{corollary}

\theoremstyle{definition}
\newaliascnt{definition}{theorem}
\newtheorem{definition}[definition]{\notesdefinitionname}
\aliascntresetthe{definition}
\newaliascnt{example}{theorem}
\newtheorem{example}[example]{\notesexamplename}
\aliascntresetthe{example}
\newaliascnt{problem}{theorem}
\newtheorem{problem}[problem]{\notesproblemname}
\aliascntresetthe{problem}
\newtheorem{exercise}{\notesexercisename}

\theoremstyle{remark}
\newaliascnt{remark}{theorem}
\newtheorem{remark}[remark]{\notesremarkname}
\aliascntresetthe{remark}

\newenvironment{solution}
{
\begin{proof}[\notessolutionname]}
  {
\end{proof}}

\usepackage[nameinlink,noabbrev]{cleveref}

\crefname{theorem}{theorem}{theorems}
\crefname{lemma}{lemma}{lemmas}
\crefname{proposition}{proposition}{propositions}
\crefname{corollary}{corollary}{corollaries}
\crefname{definition}{definition}{definitions}
\crefname{example}{example}{examples}
\crefname{problem}{problem}{problems}
\crefname{exercise}{exercise}{exercises}
\crefname{remark}{remark}{remarks}

\ifx\notelanguage\notesfinnishlanguage
\crefname{theorem}{lause}{lauseet}
\crefname{lemma}{lemma}{lemmat}
\crefname{proposition}{propositio}{propositiot}
\crefname{corollary}{seurauslause}{seurauslauseet}
\crefname{definition}{määritelmä}{määritelmät}
\crefname{example}{esimerkki}{esimerkit}
\crefname{problem}{tehtävä}{tehtävät}
\crefname{exercise}{tehtävä}{tehtävät}
\crefname{remark}{huomautus}{huomautukset}
\fi

\newcommand{\NN}{\mathbb{N}}
\newcommand{\NNone}{\mathbb{N}_1}
\newcommand{\ZZ}{\mathbb{Z}}
\newcommand{\QQ}{\mathbb{Q}}
\newcommand{\RR}{\mathbb{R}}
\newcommand{\CC}{\mathbb{C}}
\newcommand{\PP}{\mathbb{P}}

\DeclareMathOperator{\im}{im}
\DeclareMathOperator{\rank}{rank}
\DeclareMathOperator{\syt}{syt}
\DeclarePairedDelimiter{\abs}{\lvert}{\rvert}
\DeclarePairedDelimiter{\norm}{\lVert}{\rVert}
\newcommand{\set}[1]{\left\lbrace#1\right\rbrace}
\newcommand{\maxset}[1]{\max\set{#1}}
\newcommand{\minset}[1]{\min\set{#1}}
\newcommand{\powerset}[1]{\mathcal{P}\!\left(#1\right)}


\title{Johdatus logiikkaan II}
\author{Joona Piirainen}
\date{Syksy 2026 -- periodi 2}

\begin{document}

\begin{abstract}
  Luentomuistiinpanot Helsingin yliopiston kurssilta Johdatus logiikkaan II
  (MAT21015). Kurssilla on kaksi kahden tunnin luentoa viikossa aikavälillä
  26.10.2026--18.12.2026.
\end{abstract}

\maketitle

\tableofcontents

\section{Kurssitiedot}

Kurssin viralliset tiedot löytyvät
\href{https://studies.helsinki.fi/kurssit/opintojakso/otm-52c850ab-c78b-491d-86ab-14a4865e2850/MAT21015?cpId=hy-lv-77}{Helsingin
yliopiston kurssisivulta}.

\section{Muistiinpanot}

\subsection{26.10.2026: Predikaattilogiikan motivaatio ja mallit}

\subsubsection{Propositiologiikasta predikaattilogiikkaan}

Tarkastellaan päättelyä:
\begin{quote}
  Kaikki ihmiset ovat kuolevaisia. Sokrates on ihminen.
  Siis Sokrates on kuolevainen.
\end{quote}
Jos nämä kolme lausetta nimetään propositiosymboleiksi $P_0$, $P_1$ ja
$P_2$, niiden sisäinen rakenne katoaa. Propositiologiikassa
\[
  \{P_0,P_1\}\nvdash P_2:
\]
arvotus, jossa $P_0$ ja $P_1$ ovat tosia mutta $P_2$ epätosi, toteuttaa
premissit mutta ei johtopäätöstä. Propositiosymbolit eivät itsessään
ilmaise, miten ihmisyys ja kuolevaisuus liittyvät toisiinsa.

Predikaattilogiikassa puhutaan yksilöistä, niiden ominaisuuksista ja
niiden välisistä suhteista. Propositiologiikan konnektiivien lisäksi
käytetään kvanttoreita:
\begin{itemize}
  \item $\forall x$: kaikilla $x$;
  \item $\exists x$: on olemassa $x$.
\end{itemize}
Kun $I(x)$ tarkoittaa, että $x$ on ihminen, $K(x)$ tarkoittaa, että $x$
on kuolevainen, ja $s$ nimeää Sokrateen, päättelyn rakenne on
\[
  \forall x\bigl(I(x)\to K(x)\bigr),\quad I(s)
  \quad\Longrightarrow\quad K(s).
\]
Yleinen väite koskee myös Sokratesta: $I(s)\to K(s)$. Yhdessä väitteen
$I(s)$ kanssa tämä antaa johtopäätöksen $K(s)$.

\newpage
Myös matematiikassa tarvitaan tätä ilmaisukykyä. Esimerkiksi luonnollisista
luvuista voidaan tarkastella järjestystä, laskutoimituksia ja jaollisuutta.
Yleisesti tarkasteltavaan joukkoon liitetään
\begin{itemize}
  \item osajoukkoja eli yksilöiden ominaisuuksia;
  \item relaatioita, esimerkiksi järjestys- tai ekvivalenssirelaatio;
  \item funktioita, esimerkiksi yhteen- ja kertolasku;
  \item nimettyjä alkioita eli vakioita.
\end{itemize}

\subsubsection{Aakkosto: mitä symboleja käytetään?}

\begin{definition}[Aakkosto]
  Aakkosto $L$ on symbolikokoelma, jonka symbolit ovat
  \emph{relaatiosymboleja}, \emph{funktiosymboleja} tai
  \emph{vakiosymboleja}. Jokaiseen relaatio- ja funktiosymboliin liittyy
  tieto sen \emph{paikkaluvusta}, eli argumenttien lukumäärästä.
\end{definition}

Merkinnöissä $R_i^k$ ja $f_i^k$ yläindeksi $k$ ilmaisee paikkaluvun ja
alaindeksi $i$ erottaa symbolit toisistaan. Esimerkiksi $R_0^2,R_1^2,R_2^2$
ovat kaksipaikkaisia relaatiosymboleja, $f_0^2,f_1^2$ kaksipaikkaisia
funktiosymboleja ja $c_0,c_1$ vakiosymboleja. Yläindeksi ei tässä tarkoita
potenssia.

\begin{remark}[Symboli ja sen tulkinta]
  Aakkosto kertoo vain, millaisia symboleja on käytössä. Se ei vielä
  kerro, mitä ne tarkoittavat. Esimerkiksi kaksipaikkainen relaatiosymboli
  voidaan tulkita yhdessä mallissa relaatioksi $<$ ja toisessa
  relaatioksi $\geq$.
\end{remark}

\subsubsection{Malli: missä ja miten symbolit tulkitaan?}

\begin{definition}[$L$-malli]
  Olkoon $L$ aakkosto. \emph{$L$-malli} eli \emph{$L$-rakenne}
  $\mathcal M$ koostuu seuraavista tiedoista:
  \begin{enumerate}[label=(\roman*)]
    \item Epätyhjä joukko $M=\operatorname{dom}(\mathcal M)$, jota kutsutaan
      mallin \emph{universumiksi} tai \emph{perusjoukoksi}.
    \item Jokaisen vakiosymbolin $c_i\in L$ tulkinta on alkio
      \[
        c_i^{\mathcal M}\in M.
      \]
    \item Jokaisen $k$-paikkaisen relaatiosymbolin $R_i^k\in L$ tulkinta
      on relaatio
      \[
        (R_i^k)^{\mathcal M}\subseteq M^k.
      \]
    \item Jokaisen $k$-paikkaisen funktiosymbolin $f_i^k\in L$ tulkinta
      on funktio
      \[
        (f_i^k)^{\mathcal M}\colon M^k\longrightarrow M.
      \]
  \end{enumerate}
\end{definition}

Tässä $M^k$ on kaikkien $M$:n alkioista muodostettujen järjestettyjen
$k$-jonojen joukko. Erityisesti yksipaikkainen relaatio on $M$:n
osajoukko. Funktiosymbolin tulkinnan on annettava arvo jokaiselle
$M^k$:n alkiolle, ja arvon on kuuluttava samaan universumiin $M$.

\begin{example}[Reaalilukujen rakenne]
  Tarkastellaan rakennetta
  \[
    \mathcal R=(\RR,+,\cdot,<,0,1)
  \]
  aakkostossa $L=\{f_0^2,f_1^2,R_0^2,c_0,c_1\}$. Sen universumi on
  $\operatorname{dom}(\mathcal R)=\RR$, ja symbolien tulkinnat ovat
  \[
    (f_0^2)^{\mathcal R}=+,\qquad
    (f_1^2)^{\mathcal R}=\cdot,\qquad
    (R_0^2)^{\mathcal R}=<,\qquad
    c_0^{\mathcal R}=0,\quad c_1^{\mathcal R}=1.
  \]
  Yhteen- ja kertolasku ovat funktioita $\RR^2\to\RR$,
  järjestysrelaatio on $\RR^2$:n osajoukko ja vakiot ovat reaalilukuja.
\end{example}

\newpage
\begin{example}[Merkkijonoista muodostettu malli]
  Olkoon $L=\{c_0,f_0^3,R_0^1\}$ ja
  \[
    M=\{\text{miu},\text{mau}\}^{<\omega}.
  \]
  Tämä on kaikkien sanoista \text{miu} ja \text{mau} muodostettujen
  äärellisten jonojen joukko. Mukana on myös tyhjä jono $\varepsilon$.
  Esimerkiksi \text{miu}, \text{mau} ja \text{miumau} kuuluvat joukkoon.
  Määritellään tulkinnat
  \begin{align*}
    c_0^{\mathcal M}&=\varepsilon,\\
    (f_0^3)^{\mathcal M}(a,b,c)&=\varepsilon
    &&\text{kaikilla }a,b,c\in M,\\
    (R_0^1)^{\mathcal M}&=\{w\in M : w\text{ alkaa sanalla miu}\}.
  \end{align*}
  Näin saadaan $L$-malli: vakion tulkinta kuuluu universumiin,
  kolmipaikkainen funktio kuvaa $M^3$:n joukkoon $M$ ja yksipaikkainen
  relaatio on $M$:n osajoukko. Tyhjä jono $\varepsilon$ ei ole sama asia
  kuin tyhjä joukko $\varnothing$.
\end{example}

\subsubsection{Merkintätavat ja lisää esimerkkejä}

Virallisten symbolien $R_i^k,f_i^k,c_i$ sijasta voidaan käyttää
lyhyempiä nimiä, kuten $R,R',f,g,h,c,d$, kun symbolien tyypit ja
paikkaluvut ovat selvillä. Seuraavassa taulukossa mallin osat on
kirjoitettu samassa järjestyksessä kuin aakkoston symbolit on lueteltu.

\begin{center}
  \renewcommand{\arraystretch}{1.35}
  \begin{tabular}{lll}
    \hline
    Aakkosto $L$ & Tulkittavat symbolit & Esimerkki $L$-mallista\\
    \hline
    $\{c_0\}$ & vakio & $(\RR,\pi)$\\
    $\{R_0^1,R_0^2\}$ & 1- ja 2-paikkainen relaatio & $(\ZZ,\NN,<)$\\
    $\{f_5^1,c_3\}$ & 1-paikkainen funktio ja vakio & $(\ZZ,s,42)$\\
    $\{c_0,R_1^2\}$ & vakio ja 2-paikkainen relaatio & $(\ZZ,69,\geq)$\\
    $\varnothing$ & ei tulkittavia symboleja & $\CC$\\
    \hline
  \end{tabular}
\end{center}

Kolmannella rivillä $s\colon\ZZ\to\ZZ$ on seuraajafunktio $s(n)=n+1$.
Toisella rivillä $\NN$ tulkitsee yksipaikkaisen relaatiosymbolin:
ominaisuus on siis luonnollisiin lukuihin kuuluminen, vaikka universumi
on koko $\ZZ$. Tyhjän aakkoston malli on pelkkä epätyhjä joukko.

\subsubsection{Täydentävä esimerkki: jaollisuus ja kvanttorit}

Erillisessä valokuvassa oleva jaollisuustodistus havainnollistaa,
miten matematiikan päättelyissä käytetään ominaisuuksia ja kvanttoreita.
Kokonaisluvuilla jaollisuus määritellään ehdolla
\[
  a\mid b\quad\Longleftrightarrow\quad
  \exists d\in\ZZ\;(b=ad).
\]

\begin{proposition}
  Jos $a,b,c\in\ZZ$, $a\mid b$ ja $a\mid c$, niin kaikilla
  $x,y\in\ZZ$ pätee $a\mid(xb+yc)$.
\end{proposition}

\begin{proof}
  Oletusten nojalla on olemassa $d,d'\in\ZZ$, joille $b=ad$ ja $c=ad'$.
  Olkoot $x,y\in\ZZ$ mielivaltaisia. Tällöin
  \[
    xb+yc=xad+yad'=a(xd+yd').
  \]
  Koska $xd+yd'\in\ZZ$, jaollisuuden määritelmä antaa $a\mid(xb+yc)$.
\end{proof}

Todistuksessa ilmaukset \emph{on olemassa} ja \emph{mielivaltaiset}
vastaavat eksistenssi- ja universaalikvanttoria. Kun universumiksi
valitaan $\ZZ$ ja $D(a,b)$ tarkoittaa relaatiota $a\mid b$, väitteen
kvanttorirakenne on
\[
  \forall a\,\forall b\,\forall c\,
  \Bigl((D(a,b)\land D(a,c))\to
  \forall x\,\forall y\,D(a,xb+yc)\Bigr).
\]

\clearpage
\subsection{29.10.2026: Termit, atomikaavat ja kaavojen muodostaminen}

Predikaattilogiikan kaavat rakennetaan vaiheittain: ensin valitaan
\emph{termit}, sitten muodostetaan niistä \emph{atomikaavoja} ja lopuksi
liitetään kaavoja yhteen konnektiiveilla ja kvanttoreilla.
Tärkein ero on tämä: termillä viitataan yksilöön, kun taas kaavalla
ilmaistaan yksilöiden ominaisuus tai niiden välinen suhde.

\subsubsection{Kaavojen symbolit ja termit}

Aakkosto $L$ määrää käytettävissä olevat relaatio-, vakio- ja
funktiosymbolit sekä relaatioiden ja funktioiden paikkaluvut.
Niiden lisäksi käytössä ovat aina
\begin{itemize}
  \item muuttujat $x_0,x_1,x_2,\ldots$ (lyhyemmin myös $x,y,z,\ldots$);
  \item yhtäsuuruussymboli $=$ ja sulkeet;
  \item konnektiivit $\neg,\land,\lor,\to,\leftrightarrow$;
  \item kvanttorit $\forall$ ja $\exists$.
\end{itemize}
Muuttujia ja loogisia symboleja ei siis tarvitse erikseen lisätä
aakkostoon $L$. Esimerkiksi yhtäsuuruutta voi käyttää myös silloin,
kun $L=\varnothing$.

\begin{definition}[Termit ilman funktiosymboleita]
  Oletetaan tässä luennossa, ettei $L$ sisällä funktiosymboleita.
  Tällöin \emph{$L$-termejä} ovat täsmälleen kaikki muuttujat ja
  aakkostoon $L$ kuuluvat vakiosymbolit.
\end{definition}

\begin{example}[Mikä on termi?]
  Olkoon $L=\{P,R,c\}$, missä $P$ on yksipaikkainen ja $R$
  kaksipaikkainen relaatiosymboli ja $c$ on vakiosymboli.
  Tällöin $x,y,x_{17}$ ja $c$ ovat termejä. Sen sijaan $P$ ja $R$
  eivät ole termejä, eikä $P(x)$ ole termi: se on atomikaava.
  Myöskään $d$ ei ole tämän aakkoston vakio eikä $x+c$ ole sen termi,
  koska symboleita $d$ ja $+$ ei ole annettu aakkostossa.
\end{example}

Numeromerkintä, kuten $3$, ei automaattisesti kuulu muodolliseen kieleen:
sitä voidaan käyttää vakiosymbolina, jos niin on sovittu. Vakiosymboli
$c$ on puolestaan nimi, jonka tarkoittama alkio kiinnitetään mallissa.
Muuttujan arvo voi vaihdella samassakin mallissa.

\subsubsection{Atomikaavat: yksinkertaisimmat väitteet}

\begin{definition}[Atomikaavat]
  Aakkoston $L$ atomikaavat ovat seuraavaa kahta muotoa:
  \begin{enumerate}
    \item $t=u$, missä $t$ ja $u$ ovat $L$-termejä;
    \item $S(t_1,\ldots,t_n)$, missä $S\in L$ on $n$-paikkainen
      relaatiosymboli ja jokainen $t_i$ on $L$-termi.
  \end{enumerate}
\end{definition}

Atomikaava ei vielä sisällä konnektiiveja eikä kvanttoreita.
Relaation argumenttien määrän täytyy vastata sen paikkalukua.
Argumenttien järjestys on olennainen, ja sama termi saa esiintyä
useassa argumenttipaikassa.

\begin{example}[Atomikaavojen lukeminen]
  Aakkostossa $L=\{P,R,c\}$ edellä sovituilla paikkaluvuilla
  atomikaavoja ovat esimerkiksi
  \[
    x=y,\qquad x=c,\qquad P(x),\qquad P(c),\qquad
    R(x,c),\qquad R(c,c).
  \]
  Jos havainnollistavassa tulkinnassa $P$ tarkoittaa opiskelijuutta,
  $R(x,y)$ tarkoittaa, että $x$ tuntee $y$:n, ja $c$ nimeää Ainon,
  niin $P(c)$ ilmaisee Ainon olevan opiskelija ja $R(x,c)$ ilmaisee
  $x$:n tuntevan Ainon. Kaava $R(c,x)$ ilmaisee eri asian:
  Aino tuntee $x$:n. Pelkkä syntaksi ei takaa näiden väitteiden
  samanaikaista totuutta.
\end{example}

\clearpage
\subsubsection{Kaavojen induktiivinen määritelmä}

\begin{definition}[$L$-kaavat]
  Kaavat muodostetaan seuraavilla säännöillä:
  \begin{enumerate}
    \item Jokainen $L$-atomikaava on $L$-kaava.
    \item Jos $A$ ja $B$ ovat $L$-kaavoja, niin myös
      \[
        \neg A,\quad (A\land B),\quad (A\lor B),\quad
        (A\to B),\quad (A\leftrightarrow B)
      \]
      ovat $L$-kaavoja.
    \item Jos $A$ on $L$-kaava ja $x$ on muuttuja, niin
      $\forall x A$ ja $\exists x A$ ovat $L$-kaavoja.
  \end{enumerate}
  Muita kaavoja ei ole: jokainen kaava syntyy näistä säännöistä
  äärellisen monella askeleella.
\end{definition}

Kirjaimet $A$ ja $B$ ovat tässä kokonaisia kaavoja tarkoittavia
apumerkintöjä, eivät uusia aakkoston symboleja. Negaatio tarvitsee yhden
kaavan, muut konnektiivit kaksi. Kvanttori tarvitsee muuttujan ja kaavan.
Muodostussääntö ei vaadi, että kvantifioitava muuttuja esiintyisi kaavassa:
myös $\forall x P(c)$ on oikein muodostettu kaava.

\begin{example}[Kaavan rakentaminen askel askeleelta]
  Rakennetaan aakkostossa $L=\{P,R,c\}$ kaava
  \[
    \forall x\bigl(P(x)\to\exists y\,(R(x,y)\land P(y))\bigr).
  \]
  \begin{enumerate}
    \item $x$ ja $y$ ovat termejä.
    \item $P(x)$, $R(x,y)$ ja $P(y)$ ovat atomikaavoja.
    \item $(R(x,y)\land P(y))$ on kaava konjunktiosäännöllä.
    \item $\exists y\,(R(x,y)\land P(y))$ on kaava
      eksistenssikvanttorin muodostussäännöllä.
    \item Yhdistetään $P(x)$ ja edellinen kaava implikaatiolla.
    \item Lisätään koko implikaation eteen $\forall x$.
  \end{enumerate}
  Edellisen esimerkin tulkinnassa kaava sanoo:
  \emph{jokainen opiskelija tuntee jonkun opiskelijan}.
  Tunnettu opiskelija voi riippua siitä, ketä $x$ tarkoittaa.
  Kaava ei vaadi, että $x$ ja $y$ olisivat eri yksilöitä.
\end{example}

\clearpage
\subsubsection{Kvanttorien ja konnektiivien lukeminen}

Konnektiivien merkitykset ovat tutut: $\neg A$ tarkoittaa ``ei $A$'',
$A\land B$ tarkoittaa ``$A$ ja $B$'', $A\lor B$ tarkoittaa
``$A$ tai $B$ (tai molemmat)'', $A\to B$ tarkoittaa ``jos $A$, niin $B$''
ja $A\leftrightarrow B$ tarkoittaa ``$A$ jos ja vain jos $B$''.
Kvanttoreilla puhutaan tarkasteltavan universumin yksilöistä:
\[
  \forall x A:\ \text{jokaisella }x\text{ pätee }A,
  \qquad
  \exists x A:\ \text{jollakin }x\text{ pätee }A.
\]
Eksistenssikvanttori sanoo \emph{ainakin yksi}, ei \emph{täsmälleen yksi}.
Se ei myöskään tarkoita automaattisesti ``joku toinen''.

\begin{example}[Kaikki ja jokin: kaksi perusmallia]
  Olkoon $L=\{P,Q\}$, missä molemmat symbolit ovat yksipaikkaisia.
  Kun $P$ tarkoittaa opiskelijuutta ja $Q$ pyöräilyä, saadaan
  \begin{align*}
    \forall x\,(P(x)\to Q(x))
    &\quad\text{Jokainen opiskelija pyöräilee.}\\
    \exists x\,(P(x)\land Q(x))
    &\quad\text{Joku opiskelija pyöräilee.}
  \end{align*}
  Ensimmäisessä kaavassa implikaatio rajaa vaatimuksen opiskelijoihin.
  Kaava $\forall x\,(P(x)\land Q(x))$ vaatisi sen sijaan, että
  universumin \emph{jokainen} yksilö on sekä opiskelija että pyöräilijä.

  Toisessa kaavassa konjunktio vaatii samalta yksilöltä molemmat
  ominaisuudet. Kaava $\exists x\,(P(x)\to Q(x))$ ei ilmaise samaa:
  jo yksi yksilö, joka ei ole opiskelija, riittää tekemään sen todeksi.
\end{example}

\begin{example}[Sama todistaja vai eri todistajat?]
  Kun $R(x,y)$ tarkoittaa, että $x$ tuntee $y$:n, kaavat
  \begin{align*}
    \forall x\,\exists y\,R(x,y)
    &\quad\text{Jokainen tuntee jonkun.}\\
    \exists y\,\forall x\,R(x,y)
    &\quad\text{On joku, jonka kaikki tuntevat.}
  \end{align*}
  sanovat eri asioita. Jos universumissa on vain Aino ja Eero ja
  kumpikin tuntee vain itsensä, ensimmäinen väite on tosi: valitaan
  kullekin $x$:lle $y=x$. Toinen on epätosi, koska ei ole yhtä
  henkilöä, jonka molemmat tuntevat. Kvanttorien järjestystä ei siis
  yleensä voi vaihtaa.
\end{example}

\subsubsection{Sulkeet ja kaavan rakenteen purkaminen}

Kvanttorin \emph{vaikutusala} on se kaava, johon se muodostussäännössä
liitetään. Sulkeet kertovat, mitkä osat kuuluvat yhteen. Esimerkiksi
\[
  \forall x\,(P(x)\lor Q(x))
  \qquad\text{ja}\qquad
  (\forall x\,P(x))\lor Q(x)
\]
ovat eri kaavoja. Vasemmalla kvanttori koskee koko disjunktiota;
oikealla se koskee vain osaa $P(x)$, joten jälkimmäisen $Q(x)$:n
muuttuja ei kuulu tämän kvanttorin vaikutusalaan.
Käytännössä uloimpia sulkeita jätetään pois, kun rakenne pysyy selvänä.
Epäselvässä kohdassa sulkeet kannattaa aina kirjoittaa näkyviin.

Kaavan rakennetta voi tutkia myös \emph{ulkoa sisään} kysymällä,
mikä oli viimeinen muodostusaskel. Esimerkiksi kaavassa
\[
  \neg\exists x\,(P(x)\land R(x,c))
\]
uloin merkki on negaatio. Sen alla on eksistenssikvanttori ja tämän
alla konjunktio, jonka osat $P(x)$ ja $R(x,c)$ ovat atomikaavoja.
Kaava sanoo esimerkkitulkinnassa, ettei kukaan opiskelija tunne Ainoa.
Sen muodostaminen ja purkaminen käyttävät samoja sääntöjä eri suuntiin.

\clearpage
\subsubsection{Syntaksi ei ole sama asia kuin totuus}

\emph{Syntaksi} kertoo, onko merkkijono oikein muodostettu kaava.
\emph{Semantiikka} kertoo, mitä kaava ilmaisee annetussa mallissa
ja milloin se on tosi. Oikein muodostettu kaava voi olla epätosi.
Esimerkiksi $\neg(x=x)$ on täysin kelvollinen kaava, vaikka se on
aina epätosi, kun yhtäsuuruus tulkitaan tavalliseen tapaan.

\begin{example}[Tavallisten syntaksivirheiden tunnistaminen]
  Käytetään jälleen aakkostoa $L=\{P,R,c\}$, jossa $P$ on
  yksipaikkainen, $R$ kaksipaikkainen ja $c$ vakio.
  \begin{center}
    \renewcommand{\arraystretch}{1.25}
    \begin{tabular}{lp{0.67\textwidth}}
      \hline
      Merkkijono & Miksi se ei ole tämän aakkoston kaava?\\
      \hline
      $R(x)$ & Kaksipaikkainen relaatio tarvitsee kaksi termiä.\\
      $P(x,y)$ & Yksipaikkaiselle relaatiolle on annettu kaksi termiä.\\
      $P(R(x,y))$ & Relaation argumentiksi tarvitaan termi, ei kaavaa.\\
      $x\land P(x)$ & Konnektiivin molemmin puolin tarvitaan kaava.\\
      $\forall c\,P(c)$ & Kvanttorilla sidotaan muuttuja, ei vakiosymbolia.\\
      $Q(x)$ & Relaatiosymboli $Q$ ei kuulu annettuun aakkostoon.\\
      \hline
    \end{tabular}
  \end{center}
  Sen sijaan $R(x,x)$ ja $\forall x P(c)$ ovat kaavoja.
  Samojen termien toistuminen ja kvanttorin kannalta tarpeeton muuttuja
  eivät riko muodostussääntöjä.
\end{example}

\subsubsection{Kertaus ja oma tarkistus}

Kun tarkistat kaavaa, etene seuraavasti:
\begin{enumerate}
  \item Tarkista aakkosto, symbolien tyypit ja paikkaluvut.
  \item Erota yksilöihin viittaavat termit väitteitä ilmaisevista kaavoista.
  \item Tarkista atomikaavat: yhtäsuuruuden molemmin puolin ja
    relaation jokaisessa argumenttipaikassa on oltava termi.
  \item Tunnista uloin konnektiivi tai kvanttori ja tarkista sen osat
    samalla tavalla, kunnes pääset atomikaavoihin.
  \item Lue vasta sitten kaava sanallisesti sovitussa tulkinnassa.
\end{enumerate}

\begin{exercise}[Tunnista rakenne]
  Olkoon $L=\{P,R,c\}$ kuten edellä. Luokittele seuraavat termeiksi,
  atomikaavoiksi, muiksi kaavoiksi tai virheellisiksi ilmauksiksi:
  \[
    c,\quad R(c,x),\quad \neg P(c),\quad
    \exists x\,(R(c,x)\land\neg(x=c)),\quad R(P(c),x).
  \]
  Mitä neljäs ilmaus sanoo, jos $R$ tarkoittaa tuntemista ja $c$ Ainoa?
\end{exercise}

\begin{solution}
  $c$ on termi ja $R(c,x)$ atomikaava. $\neg P(c)$ on negaatiolla
  muodostettu kaava. Neljäs ilmaus on eksistenssikvanttorilla
  muodostettu kaava: \emph{Aino tuntee jonkun muun kuin itsensä}.
  Ehto $\neg(x=c)$ vaatii eri yksilön; sitä lyhennetään usein
  merkinnällä $x\ne c$. Viimeinen ilmaus on virheellinen, koska
  $P(c)$ on kaava eikä relaatioon argumentiksi kelpaava termi.
\end{solution}

\end{document}
